(1) 取长度 \(dx\) 弦元的横向动量平衡，得 \(\mu\,\partial_t^2 y=\partial_x[T(x)\partial_x y]\)。局部常系数近似的相速度是 \(c(x)=\sqrt{T(x)/\mu}=\sqrt{T_0/\mu}(1+\alpha x/L)\)，因此局部波数随 \(x\) 变化。这里 \(T'(x)y_x\) 不能在精确方程中预先删去，只能在高频近似中控制其相对阶数。
(2) 令 \(y\sim A(x)\cos[\omega t-\phi(x)]\)，最低阶WKB给 \(k(x)=\phi'(x)=\omega/c(x)\)。故 \(\phi(x)-\phi(0)=[\omega L/(\alpha c_0)]\ln(1+\alpha x/L)\)，\(c_0=\sqrt{T_0/\mu}\)。波包群时延由 \(t_{\rm travel}=\int_0^L dx/c(x)=[L/(\alpha c_0)]\ln(1+\alpha)\) 给出；\(\alpha\) 可正可负但须保持张力正值。
(3) 高频横波一个周期的平均能流为 \(\overline P=\frac12T(x)\omega k(x)A(x)^2=\frac12\omega^2\sqrt{\mu T(x)}A(x)^2\)。在无耗散、弱反射近似下它独立于 \(x\)，故 \(A(x)/A(0)=[T_0/T(x)]^{1/4}=(1+\alpha x/L)^{-1/2}\)。振幅随张力增大而减小，并不是因耗散而减少，能流仍守恒。
(4) WKB要求一个局部波长内参数变化很小，即 \(|k'|/k^2=|c'(x)|/\omega=|\alpha c_0|/(\omega L)\ll1\)，对当前线性 \(c(x)\) 是常数。用 \(\ln(1+\alpha)/\alpha\to1\)，可得 \(t_{\rm travel}\to L/c_0\)、\(A\to A(0)\)，恢复匀张力弦。若变化尺度与波长相当，左右行波耦合并出现反射，单向能流近似不再闭合。
